\documentclass[10pt,a4paper]{amsart}
\usepackage{amsmath,amssymb,amsthm,graphicx}
% Retained preamble for the Stacks Project chapter {CH} (shared project preamble).
% The chapter file pinned-original-source.tex only carries \input{{preamble}} on its first line,
% so no preamble text lives inside the chapter itself.  The Stacks Project uses one single shared
% preamble for every chapter: the snapshot stages it three times (sources/stacks-chow/preamble.tex,
% sources/stacks-etale-cohomology/preamble.tex and focused-seeds/AGP/sources/stacks-snapshot-20261011/
% preamble.tex) and those copies are identical in size.  The block below is copied mechanically,
% line for line, from sources/stacks-chow/preamble.tex lines 14-18 and 154-217: the Xy-pic setup and
% the native theorem declarations together with the mathematical macros.  Nothing was added, removed
% or redefined; the wrapper class and the amsmath/amssymb/amsthm/graphicx base packages are supplied
% by the assembler.
% For commutative diagrams we use Xy-pic
% Chapter: stacks-cohomology.tex -- Cohomology of Sheaves.
\usepackage[all]{xy}

% We use 2cell for 2-commutative diagrams.
\xyoption{2cell}

% Theorem environments.
%
\theoremstyle{plain}
\newtheorem{theorem}[subsection]{Theorem}
\newtheorem{proposition}[subsection]{Proposition}
\newtheorem{lemma}[subsection]{Lemma}

\theoremstyle{definition}
\newtheorem{definition}[subsection]{Definition}
\newtheorem{example}[subsection]{Example}
\newtheorem{exercise}[subsection]{Exercise}
\newtheorem{situation}[subsection]{Situation}

\theoremstyle{remark}
\newtheorem{remark}[subsection]{Remark}
\newtheorem{remarks}[subsection]{Remarks}

\numberwithin{equation}{subsection}

% Macros
%
\def\lim{\mathop{\mathrm{lim}}\nolimits}
\def\colim{\mathop{\mathrm{colim}}\nolimits}
\def\Spec{\mathop{\mathrm{Spec}}}
\def\Hom{\mathop{\mathrm{Hom}}\nolimits}
\def\Ext{\mathop{\mathrm{Ext}}\nolimits}
\def\SheafHom{\mathop{\mathcal{H}\!\mathit{om}}\nolimits}
\def\SheafExt{\mathop{\mathcal{E}\!\mathit{xt}}\nolimits}
\def\Sch{\mathit{Sch}}
\def\Mor{\mathop{\mathrm{Mor}}\nolimits}
\def\Ob{\mathop{\mathrm{Ob}}\nolimits}
\def\Sh{\mathop{\mathit{Sh}}\nolimits}
\def\NL{\mathop{N\!L}\nolimits}
\def\CH{\mathop{\mathrm{CH}}\nolimits}
\def\proetale{{pro\text{-}\acute{e}tale}}
\def\etale{{\acute{e}tale}}
\def\QCoh{\mathit{QCoh}}
\def\Ker{\mathop{\mathrm{Ker}}}
\def\Im{\mathop{\mathrm{Im}}}
\def\Coker{\mathop{\mathrm{Coker}}}
\def\Coim{\mathop{\mathrm{Coim}}}

% Boxtimes
%
\DeclareMathSymbol{\boxtimes}{\mathbin}{AMSa}{"02}

%
% Macros for moduli stacks/spaces
%
\def\QCohstack{\mathcal{QC}\!\mathit{oh}}
\def\Cohstack{\mathcal{C}\!\mathit{oh}}
\def\Spacesstack{\mathcal{S}\!\mathit{paces}}
\def\Quotfunctor{\mathrm{Quot}}
\def\Hilbfunctor{\mathrm{Hilb}}
\def\Curvesstack{\mathcal{C}\!\mathit{urves}}
\def\Polarizedstack{\mathcal{P}\!\mathit{olarized}}
\def\Complexesstack{\mathcal{C}\!\mathit{omplexes}}
% \Pic is the operator that assigns to X its picard group, usage \Pic(X)
% \Picardstack_{X/B} denotes the Picard stack of X over B
% \Picardfunctor_{X/B} denotes the Picard functor of X over B
\def\Pic{\mathop{\mathrm{Pic}}\nolimits}
\def\Picardstack{\mathcal{P}\!\mathit{ic}}
\def\Picardfunctor{\mathrm{Pic}}
\def\Deformationcategory{\mathcal{D}\!\mathit{ef}}

% --- appended 2026-10-12 (Chief): the project own definitions of the reference/slogan/history
% display environments, copied mechanically from the shared Stacks preamble lines 8-12
% (\usepackage{verbatim} plus the three \newenvironment lines, which the chapter preamble keeps
% verbatim from the project preamble). Nothing added, removed or redefined: these environments
% carry no mathematical content and the project itself typesets them as comments. Without these
% three lines a counted statement whose own text uses them fails with Environment slogan undefined.
\usepackage{verbatim}
\newenvironment{reference}{\comment}{\endcomment}
\newenvironment{slogan}{\comment}{\endcomment}
\newenvironment{history}{\comment}{\endcomment}
\usepackage{cleveref}
\crefname{originallemma}{original lemma}{original lemmas}
\Crefname{originallemma}{Original lemma}{Original lemmas}
\makeatletter
\newlabel{lemma-base-change-map-flat-case}{{02N7}{source}{}{sourcepointer-1}{}}
\newlabel{lemma-base-change-map-flat-case@cref}{{[originallemma][0][]02N7}{[1][0][]source}{}{}{}}
\makeatother
\crefname{originallemma}{original lemma}{original lemmas}
\Crefname{originallemma}{Original lemma}{Original lemmas}
\makeatletter
\newlabel{lemma-proper-base-change}{{09V5}{source}{}{sourcepointer-2}{}}
\newlabel{lemma-proper-base-change@cref}{{[originallemma][0][]09V5}{[1][0][]source}{}{}{}}
\makeatother
\crefname{originallemma}{original lemma}{original lemmas}
\Crefname{originallemma}{Original lemma}{Original lemmas}
\makeatletter
\newlabel{topology-lemma-base-change-separated}{{0CY4}{source}{}{sourcepointer-3}{}}
\newlabel{topology-lemma-base-change-separated@cref}{{[originallemma][0][]0CY4}{[1][0][]source}{}{}{}}
\makeatother
\crefname{originaltheorem}{original theorem}{original theorems}
\Crefname{originaltheorem}{Original theorem}{Original theorems}
\makeatletter
\newlabel{topology-theorem-characterize-proper}{{005R}{source}{}{sourcepointer-4}{}}
\newlabel{topology-theorem-characterize-proper@cref}{{[originaltheorem][0][]005R}{[1][0][]source}{}{}{}}
\makeatother
\title{Proper base change for the derived pushforward of a complex on topological spaces}
\author{The Stacks Project Authors}
\date{stacks-project:cohomology.tex:local-snapshot-20261011; GNU FDL 1.2 or later}
\begin{document}
\maketitle
\noindent\textit{Editorial note.} All delivered mathematical text is the author's own, copied from the pinned Stacks Project chapter file cohomology.tex; no statement or proof step was rewritten, reordered or removed.\par\medskip
\noindent\textit{Editorial note.} This card was re-admitted after a gate audit: it had been excluded only because the chapter's retained preamble previously lacked the reference/slogan/history environments. The shared preamble now carries them (appended 2026-10-12), so the card uses that file as its preamble with no per-card copy.\par\medskip
\noindent\textit{Editorial note.} The author's own bibliographic remark, which sits in a reference environment at the head of the theorem and points at the original exposition, is preserved in the source span and is not typeset because the project's own preamble defines reference environments as comments; no bibliography excerpt is needed.\par\medskip
\noindent\textit{Editorial note.} The label theorem-proper-base-change is defined only in cohomology.tex, so no cross-chapter restatement had to be weighed for this card.\par\medskip
\noindent\textit{Allowed standard prerequisites (editorial).}\begin{itemize}
\item Fibre formula for stalks of the derived pushforward under closedness, separatedness and quasi-compact fibres. Same chapter, Lemma lemma-proper-base-change, invoked at proof lines 3128-3135. Supplies the identification of both sides with derived global sections over the fibres.
\item Characterization of properness and stability of the relevant hypotheses under base change. Topology, Theorem topology-theorem-characterize-proper and Lemma topology-lemma-base-change-separated, invoked at proof lines 3124-3127. Supplies that the base changed map is proper with quasi-compact fibres and that both maps are closed.
\item The base change map for derived pushforwards. Same chapter, Lemma lemma-base-change-map-flat-case, named in the statement at line 3117. Supplies the map whose bijectivity is the content of the theorem.
\end{itemize}\medskip
\ifdefined\hypertarget\hypertarget{sourcepointer-1}{}\fi
\noindent\textit{Source reference.} Original lemma 02N7: sources/stacks-cohomology/pinned-original-source.tex, source lines 2977--2977. Pointer outside the counted span, invoked once for context; the load-bearing argument of the counted statement is the author's proof reproduced in full in the proof span.\par\medskip
\ifdefined\hypertarget\hypertarget{sourcepointer-2}{}\fi
\noindent\textit{Source reference.} Original lemma 09V5: sources/stacks-cohomology/pinned-original-source.tex, source lines 3059--3059. Pointer outside the counted span, invoked once for context; the load-bearing argument of the counted statement is the author's proof reproduced in full in the proof span.\par\medskip
\ifdefined\hypertarget\hypertarget{sourcepointer-3}{}\fi
\noindent\textit{Source reference.} Original lemma 0CY4: sources/stacks-chow/tags.txt, source lines 16229--16229. Pointer outside the counted span, invoked once for context; the load-bearing argument of the counted statement is the author's proof reproduced in full in the proof span.\par\medskip
\ifdefined\hypertarget\hypertarget{sourcepointer-4}{}\fi
\noindent\textit{Source reference.} Original theorem 005R: sources/stacks-chow/tags.txt, source lines 219--219. Pointer outside the counted span, invoked once for context; the load-bearing argument of the counted statement is the author's proof reproduced in full in the proof span.\par\medskip
\begin{theorem}[Proper base change]
\label{theorem-proper-base-change}
\begin{reference}
\cite[Expose V bis, 4.1.1]{SGA4}
\end{reference}
Consider a cartesian square of topological spaces
$$
\xymatrix{
X' = Y' \times_Y X \ar[d]_{f'} \ar[r]_-{g'} & X \ar[d]^f \\
Y' \ar[r]^g & Y
}
$$
Assume that $f$ is proper.
Let $E$ be an object of $D^+(X)$. Then the base change map
$$
g^{-1}Rf_*E \longrightarrow Rf'_*(g')^{-1}E
$$
of Lemma \ref{lemma-base-change-map-flat-case} is an isomorphism
in $D^+(Y')$.
\end{theorem}

\begin{proof}
Let $y' \in Y'$ be a point with image $y \in Y$. It suffices to show that
the base change map induces an isomorphism on stalks at $y'$.
As $f$ is proper it follows that $f'$ is proper, the
fibres of $f$ and $f'$ are quasi-compact and $f$ and $f'$ are closed, see
Topology, Theorem \ref{topology-theorem-characterize-proper} and
Lemma \ref{topology-lemma-base-change-separated}.
Thus we can apply Lemma \ref{lemma-proper-base-change} twice to see that
$$
(Rf'_*(g')^{-1}E)_{y'} = R\Gamma((f')^{-1}(y'), (g')^{-1}E|_{(f')^{-1}(y')})
$$
and
$$
(Rf_*E)_y = R\Gamma(f^{-1}(y), E|_{f^{-1}(y)})
$$
The induced map of fibres $(f')^{-1}(y') \to f^{-1}(y)$ is
a homeomorphism of topological spaces and the pull back of
$E|_{f^{-1}(y)}$ is $(g')^{-1}E|_{(f')^{-1}(y')}$. The
desired result follows.
\end{proof}

\end{document}
